%-----------------------------------------------------------------------------------------------------------------------------------------------%
%	The MIT License (MIT)
%
%	Copyright (c) 2019 Jan Küster
%
%	Permission is hereby granted, free of charge, to any person obtaining a copy
%	of this software and associated documentation files (the "Software"), to deal
%	in the Software without restriction, including without limitation the rights
%	to use, copy, modify, merge, publish, distribute, sublicense, and/or sell
%	copies of the Software, and to permit persons to whom the Software is
%	furnished to do so, subject to the following conditions:
%	
%	THE SOFTWARE IS PROVIDED "AS IS", WITHOUT WARRANTY OF ANY KIND, EXPRESS OR
%	IMPLIED, INCLUDING BUT NOT LIMITED TO THE WARRANTIES OF MERCHANTABILITY,
%	FITNESS FOR A PARTICULAR PURPOSE AND NONINFRINGEMENT. IN NO EVENT SHALL THE
%	AUTHORS OR COPYRIGHT HOLDERS BE LIABLE FOR ANY CLAIM, DAMAGES OR OTHER
%	LIABILITY, WHETHER IN AN ACTION OF CONTRACT, TORT OR OTHERWISE, ARISING FROM,
%	OUT OF OR IN CONNECTION WITH THE SOFTWARE OR THE USE OR OTHER DEALINGS IN
%	THE SOFTWARE.
%	
%
%-----------------------------------------------------------------------------------------------------------------------------------------------%


%============================================================================%
%
%	DOCUMENT DEFINITION
%
%============================================================================%

%we use article class because we want to fully customize the page and don't use a cv template
\documentclass[10pt,A4]{article}	


%----------------------------------------------------------------------------------------
%	ENCODING
%----------------------------------------------------------------------------------------

% we use utf8 since we want to build from any machine
\usepackage[utf8]{inputenc}		

%----------------------------------------------------------------------------------------
%	LOGIC
%----------------------------------------------------------------------------------------

% provides \isempty test
\usepackage{xstring, xifthen}

%----------------------------------------------------------------------------------------
%	FONT BASICS
%----------------------------------------------------------------------------------------

% some tex-live fonts - choose your own

%\usepackage[defaultsans]{droidsans}
%\usepackage[default]{comfortaa}
%\usepackage{cmbright}
\usepackage[default]{raleway}
%\usepackage{fetamont}
%\usepackage[default]{gillius}
%\usepackage[light,math]{iwona}
%\usepackage[thin]{roboto} 

% set font default
\renewcommand*\familydefault{\sfdefault} 	
\usepackage[T1]{fontenc}

% more font size definitions
\usepackage{moresize}

%----------------------------------------------------------------------------------------
%	FONT AWESOME ICONS
%---------------------------------------------------------------------------------------- 

% include the fontawesome icon set
\usepackage{fontawesome}

% use to vertically center content
% credits to: http://tex.stackexchange.com/questions/7219/how-to-vertically-center-two-images-next-to-each-other
\newcommand{\vcenteredinclude}[1]{\begingroup
\setbox0=\hbox{\includegraphics{#1}}%
\parbox{\wd0}{\box0}\endgroup}

% use to vertically center content
% credits to: http://tex.stackexchange.com/questions/7219/how-to-vertically-center-two-images-next-to-each-other
\newcommand*{\vcenteredhbox}[1]{\begingroup
\setbox0=\hbox{#1}\parbox{\wd0}{\box0}\endgroup}

% icon shortcut
\newcommand{\icon}[3] { 							
	\makebox(#2, #2){\textcolor{maincol}{\csname fa#1\endcsname}}
}	

% icon with text shortcut
\newcommand{\icontext}[4]{ 						
	\vcenteredhbox{\icon{#1}{#2}{#3}}  \hspace{2pt}  \parbox{0.9\mpwidth}{\textcolor{#4}{#3}}
}

% icon with website url
\newcommand{\iconhref}[5]{ 						
    \vcenteredhbox{\icon{#1}{#2}{#5}}  \hspace{2pt} \href{#4}{\textcolor{#5}{#3}}
}

% icon with email link
\newcommand{\iconemail}[5]{ 						
    \vcenteredhbox{\icon{#1}{#2}{#5}}  \hspace{2pt} \href{mailto:#4}{\textcolor{#5}{#3}}
}

%----------------------------------------------------------------------------------------
%	PAGE LAYOUT  DEFINITIONS
%----------------------------------------------------------------------------------------

% page outer frames (debug-only)
% \usepackage{showframe}		

% we use paracol to display breakable two columns
\usepackage{paracol}

% define page styles using geometry
\usepackage[a4paper]{geometry}

% remove all possible margins
\geometry{top=1cm, bottom=1cm, left=1cm, right=1cm}

\usepackage{fancyhdr}
\pagestyle{empty}

% space between header and content
% \setlength{\headheight}{0pt}

% indentation is zero
\setlength{\parindent}{0mm}

%----------------------------------------------------------------------------------------
%	TABLE /ARRAY DEFINITIONS
%---------------------------------------------------------------------------------------- 

% extended aligning of tabular cells
\usepackage{array}

% custom column right-align with fixed width
% use like p{size} but via x{size}
\newcolumntype{x}[1]{%
>{\raggedleft\hspace{0pt}}p{#1}}%


%----------------------------------------------------------------------------------------
%	GRAPHICS DEFINITIONS
%---------------------------------------------------------------------------------------- 

%for header image
\usepackage{graphicx}

% use this for floating figures
% \usepackage{wrapfig}
% \usepackage{float}
% \floatstyle{boxed} 
% \restylefloat{figure}

%for drawing graphics		
\usepackage{tikz}				
\usetikzlibrary{shapes, backgrounds,mindmap, trees}

%----------------------------------------------------------------------------------------
%	Color DEFINITIONS
%---------------------------------------------------------------------------------------- 
\usepackage{transparent}
\usepackage{color}

% primary color
\definecolor{maincol}{RGB}{ 225, 0, 0 }

% accent color, secondary
% \definecolor{accentcol}{RGB}{ 250, 150, 10 }

% dark color
\definecolor{darkcol}{RGB}{ 70, 70, 70 }

% light color
\definecolor{lightcol}{RGB}{245,245,245}


% Package for links, must be the last package used
\usepackage[hidelinks]{hyperref}

% returns minipage width minus two times \fboxsep
% to keep padding included in width calculations
% can also be used for other boxes / environments
\newcommand{\mpwidth}{\linewidth-\fboxsep-\fboxsep}
	


%============================================================================%
%
%	CV COMMANDS
%
%============================================================================%

%----------------------------------------------------------------------------------------
%	 CV LIST
%----------------------------------------------------------------------------------------

% renders a standard latex list but abstracts away the environment definition (begin/end)
\newcommand{\cvlist}[1] {
	\begin{itemize}{#1}\end{itemize}
}

%----------------------------------------------------------------------------------------
%	 CV TEXT
%----------------------------------------------------------------------------------------

% base class to wrap any text based stuff here. Renders like a paragraph.
% Allows complex commands to be passed, too.
% param 1: *any
\newcommand{\cvtext}[1] {
	\begin{tabular*}{1\mpwidth}{p{0.98\mpwidth}}
		\parbox{1\mpwidth}{#1}
	\end{tabular*}
}

%----------------------------------------------------------------------------------------
%	CV SECTION
%----------------------------------------------------------------------------------------

% Renders a a CV section headline with a nice underline in main color.
% param 1: section title
\newcommand{\cvsection}[1] {
	\vspace{14pt}
	\cvtext{
		\textbf{\LARGE{\textcolor{darkcol}{\uppercase{#1}}}}\\[-4pt]
		\textcolor{maincol}{ \rule{0.1\textwidth}{2pt} } \\
	}
}

%----------------------------------------------------------------------------------------
%	META SKILL
%----------------------------------------------------------------------------------------

% Renders a progress-bar to indicate a certain skill in percent.
% param 1: name of the skill / tech / etc.
% param 2: level (for example in years)
% param 3: percent, values range from 0 to 1
\newcommand{\cvskill}[3] {
	\begin{tabular*}{1\mpwidth}{p{0.72\mpwidth}  r}
 		\textcolor{black}{\textbf{#1}} & \textcolor{maincol}{#2}\\
	\end{tabular*}%
	
	\hspace{4pt}
	\begin{tikzpicture}[scale=1,rounded corners=2pt,very thin]
		\fill [lightcol] (0,0) rectangle (1\mpwidth, 0.15);
		\fill [maincol] (0,0) rectangle (#3\mpwidth, 0.15);
  	\end{tikzpicture}%
}


%----------------------------------------------------------------------------------------
%	 CV EVENT
%----------------------------------------------------------------------------------------

% Renders a table and a paragraph (cvtext) wrapped in a parbox (to ensure minimum content
% is glued together when a pagebreak appears).
% Additional Information can be passed in text or list form (or other environments).
% the work you did
% param 1: time-frame i.e. Sep 14 - Jan 15 etc.
% param 2:	 event name (job position etc.)
% param 3: Customer, Employer, Industry
% param 4: Short description
% param 5: work done (optional)
% param 6: technologies include (optional)
% param 7: achievements (optional)
\newcommand{\cvevent}[7] {
	
	% we wrap this part in a parbox, so title and description are not separated on a pagebreak
	% if you need more control on page breaks, remove the parbox
	\parbox{\mpwidth}{
		\begin{tabular*}{1\mpwidth}{p{0.72\mpwidth}  r}
	 		\textcolor{black}{\textbf{#2}} & \colorbox{maincol}{\makebox[0.25\mpwidth]{\textcolor{white}{#1}}} \\
			\textcolor{maincol}{\textbf{#3}} & \\
		\end{tabular*}\\[8pt]
	
		\ifthenelse{\isempty{#4}}{}{
			\cvtext{#4}\\
		}
	}

	\ifthenelse{\isempty{#5}}{}{
		\vspace{9pt}
		{#5}
	}

	\ifthenelse{\isempty{#6}}{}{
		\vspace{9pt}
		\cvtext{\textbf{Technologies include:}}\\
		{#6}
	}
%
%	\ifthenelse{\isempty{#7}}{}{
%		\vspace{9pt}
%		\cvtext{\textbf{Achievements include:}}\\
%		{#7}
%	}
	\vspace{14pt}
}

%----------------------------------------------------------------------------------------
%	 CV META EVENT
%----------------------------------------------------------------------------------------

% Renders a CV event on the sidebar
% param 1: title
% param 2: subtitle (optional)
% param 3: customer, employer, etc,. (optional)
% param 4: info text (optional)
\newcommand{\cvmetaevent}[4] {
	\textcolor{maincol} {\cvtext{\textbf{\begin{flushleft}#1\end{flushleft}}}}

	\ifthenelse{\isempty{#2}}{}{
	\textcolor{darkcol} {\cvtext{\textbf{#2}} }
	}

	\ifthenelse{\isempty{#3}}{}{
		\cvtext{{ \textcolor{darkcol} {#3} }}\\
	}

	\cvtext{#4}\\[14pt]
}

%---------------------------------------------------------------------------------------
%	QR CODE
%----------------------------------------------------------------------------------------

% Renders a qrcode image (centered, relative to the parentwidth)
% param 1: percent width, from 0 to 1
\newcommand{\cvqrcode}[1] {
	\begin{center}
		\includegraphics[width={#1}\mpwidth]{qrcode}
	\end{center}
}


%============================================================================%
%
%
%
%	DOCUMENT CONTENT
%
%
%
%============================================================================%
\begin{document}
\columnratio{0.31}
\setlength{\columnsep}{2.2em}
\setlength{\columnseprule}{4pt}
\colseprulecolor{lightcol}
\begin{paracol}{2}
\begin{leftcolumn}
%---------------------------------------------------------------------------------------
%	META IMAGE
%----------------------------------------------------------------------------------------
\includegraphics[width=\linewidth]{me.jpg}	%trimming relative to image size

%---------------------------------------------------------------------------------------
%	META SKILLS
%----------------------------------------------------------------------------------------
\cvsection{SKILLS}

\cvskill{Linux} {8+ yrs} {1} \\[-2pt]

\cvskill{AWS} {4+ yrs} {1} \\[-2pt]

\cvskill{Terraform} {4+ yrs} {1} \\[-2pt]

\cvskill{Kubernetes} {4+ yrs} {1} \\[-2pt]

\cvskill{CI/CD} {4+ yrs} {1} \\[-2pt]

\cvskill{Monitoring tools} {5+ yrs} {0.72} \\[-2pt]

\cvskill{Observability} {3+ yrs} {0.3} \\[-2pt]

\cvskill{Bash} {6+ yrs} {1} \\[-2pt]

\cvskill{Python} {3+ yrs} {0.4} \\[-2pt]

\cvskill{Golang} {2+ yrs} {0.2} \\[-2pt]



\vfill\null
\cvsection{CONTACT}
	
\icontext{MapMarker}{12}{Bogotà - Colombia}{black}\\[6pt]
\icontext{MobilePhone}{12}{+57 3138139592}{black}\\[6pt]
\iconhref{Linkedin}{12}{Brayan Bautista}{https://www.linkedin.com/in/brayanbautista/}{black}\\[6pt]
\iconhref{Twitter}{12}{@Braybaut}{https://twitter.com/braybaut}{black}\\[6pt]
%\iconhref{Medium}{12}{braybaut.dev}{https://braybaut.dev}{black}\\[6pt]
\iconhref{Github}{12}{braybaut}{https://github.com/braybaut}{black}\\[6pt]
\iconemail{Envelope}{12}{braybaut@gmail.com}{braybaut@gmail.com}{black}\\[6pt]

\vfill\null
%\cvqrcode{0.7}

%---------------------------------------------------------------------------------------
%	EDUCATION
%----------------------------------------------------------------------------------------
\newpage
\cvsection{EDUCATION}

\cvmetaevent
{2014 - 2019}
{B. System engineering}
{Universidad Cooperativa de Colombia}
%{The topic for the bachelor's thesis was 'Feshbach resonance'. A numerical application was built to calculate the diagrams.}

\cvsection{LANGUAGES}

\cvskill{English - Advanced }{} {0.8} \\[-2pt]

\cvskill{Spanish - Natie }{} {1} \\[-2pt]

%---------------------------------------------------------------------------------------
%	CERTIFICATION
%----------------------------------------------------------------------------------------
\cvsection{CERTIFICATIONS}

\cvmetaevent
{Redhat Certified System Administrator (RHCSA)}
{RedHat}
{}
{}

\cvmetaevent
{Scrum Fundamentals Certified (SFC)}
{VMEdu Inc.}
{}
{}

\cvsection{INTEREST}
{\cvlist{
\item Observability
\item Kubernetes
\item Development infra tools
\item opensource
\item Golang
}}

\vfill\null


\newpage


\vfill


\newpage
\mbox{} % hotfix to place qrcode on the bottom when there are not other elements
\vfill


\end{leftcolumn}
\begin{rightcolumn}
%---------------------------------------------------------------------------------------
%	TITLE  HEADER
%----------------------------------------------------------------------------------------
\fcolorbox{white}{darkcol}{\begin{minipage}[c][3.5cm][c]{1\mpwidth}
	\begin {center}
		\HUGE{ \textbf{ \textcolor{white}{ \uppercase{ Brayan Bautista } } } } \\[-24pt]
		\textcolor{white}{ \rule{0.1\textwidth}{1.25pt} } \\[4pt]
		\large{ \textcolor{white} {SRE and Sr DevOps Engineer} }
	\end {center}
\end{minipage}} \\[14pt]
\vspace{-12pt}

%---------------------------------------------------------------------------------------
%	PROFILE
%----------------------------------------------------------------------------------------
\vfill\null
\cvsection{PROFILE}

\cvtext{
I'm a System Engineering and am working as a Cloud/SRE and DevOps Engineer at Globant where I help
customers to design architectures to make scalable infrastructures and creating workflows to deploy in
production environments with zero downtime. \\

I'm also Container and OpenSource Enthusiast, I really enjoy writing about anything that might help other engineers out with topics i can see in my daily basis, I'm interested in anything about containers, serverless  and distributed systems, I really enjoy learning new languages and tools that can help in my career path, I also enjoy participate as a speaker at local meetups that promote the OpenSource. \\



}

%---------------------------------------------------------------------------------------
%	WORK EXPERIENCE
%----------------------------------------------------------------------------------------

\vfill\null
\cvsection{WORK EXPERIENCE}
\cvevent
	{ Jan2018 - Current }
	{Site reliability engineering and Sr DevOps engineer }
	{{\href{https://www.globant.com/}{Globant}}}
	{Working as SRE and SR DevOps engineer for 300+ microservices companies in US and Europe, designing scalable architectures, supporting multicultural teams in adopting DevOps culture and implementing SRE principles, I am part of the gatekepper team to interview the company's new hires.}
	{\cvlist{
		\item Define and implement architecture to support large volumes of logs with Fluentd, elastic search and Kibana.
		\item Design and implement architectures for microservices in AWS with ECS and EKS.
		\item Support teams in defining SLI/SLO metrics for the availability of productive services.
		\item Manage and deploy more than 10 Kubernetes clusters and more than 100 workloads in production both Kubernetes as a service and self hosted Kubernetes.
		\item define networking rules for kubernetes clusters with Istio service Mesh
		\item Develop CLI and tools to help developers adopt DevOps and adapt to SRE practices.
        \item Assist development teams in defining pipelines and infrastructure for new microservices.
        \item Define metrics and monitoring strategies for critical and productive services.
        \item Manage large critical systems with a high volume of transactions per second.
        \item led more than 8 DevOps engineers helping them in their career growth within the company. 
customers
	}}
	{\cvlist {
		\item AWS Core services: VPC, EC2, s3, Aurora, RDS, Balancers, AWS Lambda 
		\item Terraform, Ansible
		\item Datadog, NewRelic, Pagerduty 
		\item Docker, Kubernetes and control plane components
		\item Observability stack: Prometheus, Grafana, OpenTelemetry, Kibana, FluentD
		\item Golang, Python, Bash
		\item CI/CD Jenkins - Codebuild - Codedeploy
		\item Istio service mesh 
	}}
	{\cvlist{
		\item Ansible module for SLES12 System + Package registration
		\item Python tool for ACL-administration in cloud
		\item Fully automated migration of old systems into cloud with Ansible playbooks/roles
		\item Django tool on LDAP - Schema Review
	}}


\vfill\null
\cvsection{WORK EXPERIENCE}
\cvevent
	{ May2017 - Jan2018 }
	{SysAdmin Engineer Semisenior Advanced}
	{{\href{https://www.globant.com/}{Globant}}}
	{Worked as a Sysadmin engineer in the cloud side, supporting high volumes of instances and creating pipelines to deliver critical applications}
	{\cvlist{
		\item led the project of “Terraform Migration”, whole platform is maintained from
Terraform, He is responsible of Migration the Infrastructure maked from Terraform to
the latest Version
		\item led the build of infrastructure to new Microservices of a banking customer, he must
make the infrastructure necessary (ECS, DynamoDB, S3 Bucket, Load Balancer, etc
) also, he must make the CI/CD strategies for the new Microservices with Jenkins.
Design Continuous Integration
		\item Design Continuous Integration and Delivery workflow for Critical Applications with
Jenkins, Bamboo and Harness
		\item Collaborated with cross-functional development team members to analyze
potential system solutions based on evolving client requirements.
		\item Propose and create cloud architectures in AWS/GCP for customers using good
practices of infrastructure and security
        \item Creating productive environments for microservices-based applications with
orchestration technologies such as kubernetes and ECS
        \item Worked directly with agile Application Development teams to provide daily
support aligned with amodel of Continuous Integration/Continuous Delivery
        \item Managed production and critical cloud environments for Banking and Financial
customers
	}}
	{\cvlist {
		\item AWS Core services: VPC, EC2, s3, DynamoDB, RDS, Balancers 
		\item Terraform - Terratest
		\item Jenkins - Bamboo - Harness to develop critical Pipelines 
		\item ECS and Kubernetes to support containers workloads
	}}
	{\cvlist{
		\item Ansible module for SLES12 System + Package registration
		\item Python tool for ACL-administration in cloud
		\item Fully automated migration of old systems into cloud with Ansible playbooks/roles
		\item Django tool on LDAP - Schema Review
	}}

\vfill\null

\cvevent
	{2013 - 2017 }
	{Infrastructure Administrator}
	{\href{https://www.seaq.co/}{SEAQ Servicios}}
	{I worked as Infrastructure administrator assisting clients with their on-premise infrastructures with linux, working on different projects in the banking, health and education sectors.}
	{\cvlist{
	    \item I wrote scripts in Python and Bash for the automation of tasks within the infrastructure.
	    \item define maintenance window strategies for 600+ servers for academic and banking companies 
		\item Worked with customers to resolve Linux critical problems, improve operations and
provide exceptional customer service.
		\item Maintenance +200 Linux Red hat Virtual Machine for Educational, Financial and
Banking Customers
		\item Implemented virtualization environments for critical applications
		\item Implemented Jboss in Cluster for running Liferay CMS with technologies as
postgresql, Jackrabbit, at District institute of risk management and climate
change.
	}}
	{\cvlist {
		\item Nagios and Zabbix to support high infrastructure components
		\item Asterisk to support VoIP Infrastructure
		\item Python and Bash to write scripts 
		\item Linux CentOS - RedHat to be used as a OS base to install infrastructure services 
	}}
	{\cvlist{
		\item A web based tool for an intuitive role assignment and administration
		\item Online overview of company structures, projects, etc
		\item Tools for role review, reporting and troubleshooting
	}}


% hotfixes to create fake-space to ensure the whole height is used
\mbox{}
\vfill
\mbox{}
\vfill
\mbox{}
\vfill
\mbox{}
\end{rightcolumn}
\end{paracol}
\end{document}

